\begin{claim}{F.6}{Proved}{LM4-F.b: open arithmetic image for every finite geometric subgroup}
Let $2g-2+\nu>0$, let $M=\Gamma^+_{g,\nu}$ permit permutations
of the $\nu$ punctures, and retain
\[
 P=\widehat\pi_{g,\nu},\qquad
 B=\Out_{\rm lac}(P),\qquad H=\overline M\leq B.
\]
Write $A_{\rm ar}\leq B$ for the image of universal \emph{arithmetic}
\'etale monodromy of the moduli stack of smooth genus-$g$ curves
with an unordered reduced divisor of degree $\nu$, over $\Q$.
Then $H\triangleleft A_{\rm ar}$, and put
$\Gamma_{\rm ar}=A_{\rm ar}/H$.
For every finite subgroup $G\leq M$, the continuous homomorphism
\[
 C_{A_{\rm ar}}(G)\longrightarrow\Gamma_{\rm ar}
                                                        \tag{F.6.1}
\]
has open image.  Its kernel is $C_H(G)$.
In the notation $N=N_B(H)$ and $q:N\to N/H$, this is exactly
the assertion that
\[
 q(C_N(G))\cap\Gamma_{\rm ar}
       =q(C_{A_{\rm ar}}(G))
       \quad\hbox{is open in }\Gamma_{\rm ar}.           \tag{F.6.2}
\]
The same assertion holds with ordered punctures and the pure
oriented group throughout.  No identification of $H$ with the
abstract completion of $M$ is used.
\end{claim}
\begin{proof}
We first type the arithmetic group and its comparison maps.
Denote the indicated moduli stack over $\Q$ by $\mathscr M$,
and choose a geometric base point and an orientation-preserving
topological marking of its complex fiber.  Its universal
punctured curve supplies the outer monodromy homomorphism
\[
 R:\Pi_{\rm ar}=\pi_1^{\rm et}(\mathscr M)
           \longrightarrow B.
\]
The geometric comparison and the arithmetic exact sequence are
\[
 1\longrightarrow\widehat M
  \longrightarrow\Pi_{\rm ar}
  \longrightarrow G_{\Q}\longrightarrow1,
 \qquad R(\widehat M)=H.                               \tag{F.6.3}
\]
These are the constructions in \cite[\S\S7.1--7.2,
equation (29) and the universal \'etale monodromy
representation]{R-BoggiAut2026}.  For the geometric comparison,
finite \'etale covers analytify to finite topological covers;
the inverse comparison is compatible on an \'etale atlas and
its double and triple overlaps, so descends to the stack.
The marked complex moduli stack has contractible Teichm\"uller
space as its covering, with group $M$, which identifies its
topological fundamental group.  The universal curve then
identifies its monodromy with the completed geometric action
on $P$, rather than introducing a separate representation.
All these maps preserve peripheral conjugacy classes, since
the removed divisor is carried to itself in the universal
family.  Their base-point ambiguities are inner automorphisms
of $P$, already removed in $B$; a change of topological marking
conjugates the displayed subgroups by the corresponding element
of $M$.

Set $A_{\rm ar}=R(\Pi_{\rm ar})$.  It is compact and hence
closed in $B$.  Normality in (F.6.3) gives
$H\triangleleft A_{\rm ar}\leq N$.  The quotient map induces
a continuous surjection
\[
 \overline R:G_{\Q}\longrightarrow\Gamma_{\rm ar}.     \tag{F.6.4}
\]
The proof uses this quotient of $G_{\Q}$ as defined.  Its kernel
and the kernel of $\widehat M\to H$ need not be computed.

We next construct a number-field curve with the specified
\emph{labelled topological} $G$-action.  Finite realization
C.5 supplies a compact Riemann surface $X$, a reduced marked
divisor $D$, and a faithful holomorphic action of $G$ preserving
$D$ whose marked mapping classes are the given inclusion
$G\hookrightarrow M$.  Choose projective equations for $X$,
equations for $D$, and formulas for the finitely many
automorphisms and their inverses.  Their finitely many
coefficients lie in a finitely generated $\Q$-subalgebra
$T\subset\C$.  Adjoin the coefficients needed for their
identities and invert finitely many nonzero elements.  This
gives a proper smooth family of curves $\mathcal X\to V$,
a finite \'etale reduced divisor $\mathcal D\to V$, and
the same labelled $G$-action over a nonempty smooth open
subscheme $V\subset\operatorname{Spec}T$.  Smoothness and
disjointness are open conditions.  Faithfulness can be kept
on this open set as well: for each nonidentity $h\in G$,
the equalizer of $h$ and the identity is a proper closed
subset of the generic fiber, and this property persists
after shrinking.  The original inclusion $T\hookrightarrow\C$
defines a point of every such nonempty open, since it has
generic underlying point.

Base change to $\overline\Q$ and take the connected component
whose complex base change contains that point.  After a
finite extension this component and the family are defined
over a number field.  It is a smooth connected variety; its
complex points form a connected, locally path-connected
complex manifold and hence a path-connected space.  It
contains an $\overline\Q$-point, by taking a closed point
in any nonempty affine open.  Choose a piecewise smooth
complex path from the original fiber to this fiber.

For precision, parallel transport along this path preserves
the labelled action.  A smooth horizontal splitting of the
relative tangent sequence can be chosen to agree with the
tangent splitting on $\mathcal D$: locally the divisor is
a disjoint union of sections, and a partition of unity
extends their prescribed splittings.  Average this splitting
over the finite group $G$.  Averaging retains its right-inverse
property and its prescribed values on the divisor, and
makes it $G$-equivariant.  The horizontal lift along the
compact path has complete transport, because the family
is proper and its fibers are compact.  Its transport is an
orientation-preserving diffeomorphism of pairs commuting
with each labelled element of $G$.  Thus the endpoint curve
has exactly the required topological action.  This constructs
a deformation and its comparison diffeomorphism; it does not
assert that the original complex curve is isomorphic to a
curve over $\overline\Q$.

All coefficients of the endpoint curve, divisor, and the
finitely many labelled automorphisms lie in one finite
extension $K/\Q$.  Let $x_K:\operatorname{Spec}K\to\mathscr M$
be the resulting moduli point.  Choose the geometric path
to $x_K$ supplied by the transport just constructed and the
finite-cover comparison.  The morphism $x_K$ then supplies
a continuous homomorphism
\[
 s_K:G_K\longrightarrow\Pi_{\rm ar}
\]
whose composite with $\Pi_{\rm ar}\to G_{\Q}$ is the
inclusion $G_K\hookrightarrow G_{\Q}$.  Under the transported
marking, $R\circ s_K$ is the Galois outer action on the
punctured endpoint curve.

Every $h\in G$ is an automorphism over $K$ of that curve and
its marked divisor.  Naturality of the fiber functor on finite
\'etale covers gives
\[
 (R\circ s_K)(\sigma)\,h
       =h\,(R\circ s_K)(\sigma)
 \quad\hbox{in }B\qquad(\sigma\in G_K).                \tag{F.6.5}
\]
Indeed, pullback by $h$ commutes with scalar conjugation by
$\sigma$; choosing paths for the individual base points can
alter the induced automorphisms only by inner automorphisms
of $P$.  Passing to $B$ gives precisely the equality written.
The equivariant transport aligns these $h$ with the original
labelled inclusion, so (F.6.5) places $R(s_K(G_K))$ in
$C_{A_{\rm ar}}(G)$ itself.

The image of (F.6.1) therefore contains
$\overline R(G_K)$.  The latter has index at most
$[K:\Q]$ in $\Gamma_{\rm ar}$ by the surjectivity of
(F.6.4).  It is compact and hence closed; a closed finite-index
subgroup of a profinite group is open, as proved in F.3.
Every subgroup containing it is a union of its open cosets.
This proves openness in (F.6.1), with kernel $C_H(G)$ by
intersection with $H$.

Finally $q^{-1}(\Gamma_{\rm ar})=A_{\rm ar}$, since
$H\subset A_{\rm ar}$ and $\Gamma_{\rm ar}=q(A_{\rm ar})$.
An element of $C_N(G)$ whose quotient lies in $\Gamma_{\rm ar}$
therefore already belongs to $C_{A_{\rm ar}}(G)$.
This proves both inclusions in (F.6.2).  Ordering the divisor
and using the pure moduli stack changes none of the maps or
arguments when $G$ belongs to the pure oriented group.

For \S30, folios 195--196, scans 385--387, formula (F.6.2)
is the arithmetic open-image assertion with its exact
quotient made explicit.  F.3 supplies the centralizer image
in the entire quotient $N/H$; F.6 supplies its openness
after intersection with the arithmetic image for every
finite geometric $G$.  The inclusion
$\Gamma_{\rm ar}\hookrightarrow N/H$ is the comparison
map between the two statements.  Its image has not been
identified with all of $N/H$ here, so no larger openness
assertion is obtained by suppressing that inclusion.
\end{proof}
